\documentclass[a4paper,11pt]{article}

\usepackage[T1]{fontenc}
\usepackage[utf8]{inputenc}
\usepackage[brazil]{babel}
\usepackage{lmodern}
\usepackage{microtype}
\usepackage{amsmath,amssymb}
\usepackage{array,booktabs,longtable,tabularx}
\usepackage{enumitem}
\usepackage{geometry}
\usepackage{hyperref}
\usepackage{url}
\usepackage{xcolor}
\usepackage{verbatim}

\geometry{margin=2.5cm}
\hypersetup{
  colorlinks=true,
  linkcolor=blue,
  urlcolor=blue,
  citecolor=blue,
  pdftitle={Avaliação prática de AlphaTensor-Quantum em escala reduzida com nova métrica estrutural Clifford--não-Clifford e comparação contra baselines de ZX-calculus},
  pdfauthor={Plano técnico para projeto acadêmico}
}

\newcolumntype{Y}{>{\raggedright\arraybackslash}X}
\setlist[itemize]{noitemsep,topsep=3pt}
\setlist[enumerate]{noitemsep,topsep=3pt}

\title{Avaliação prática de AlphaTensor-Quantum em escala reduzida com nova métrica estrutural Clifford--não-Clifford e comparação contra baselines de ZX-calculus}
\author{Plano técnico para projeto acadêmico}
\date{\today}

\begin{document}
\maketitle
\tableofcontents
\bigskip

\section{Problem Statement}

Este projeto será executado como um \textbf{estudo experimental de otimização de circuitos quânticos com AlphaTensor-Quantum em escala reduzida, com introdução e validação de uma métrica de separação/estrutura Clifford--não-Clifford}. O objetivo não é reproduzir integralmente o treinamento do artigo-base, mas sim construir um pipeline público, verificável e executável que una compilação de circuitos, ressíntese, simplificação por ZX-calculus, uso criterioso do demo público do AlphaTensor-Quantum e análise quantitativa baseada em métricas estruturais.

A decisão central do projeto é metodológica: o repositório público do AlphaTensor-Quantum será tratado como uma \textbf{implementação pública parcial}. O componente executável de treino é o \texttt{src/demo}; o diretório \texttt{/src} não deve ser assumido como pipeline completo reproduzível de ponta a ponta; e o diretório \texttt{/decompositions} deve ser tratado como fonte de resultados públicos reaproveitáveis, não como prova de reprodução interna do treinamento do artigo. Por isso, o projeto é formalmente redefinido como \textbf{reprodução parcial + extensão metodológica + avaliação rigorosa}.

A pergunta central de pesquisa é a seguinte: \textit{uma família de métricas de estrutura Clifford--não-Clifford consegue explicar, ao menos parcialmente, o ganho de otimização em T-count obtido por diferentes métodos públicos, em especial PyZX, QuiZX+\texttt{circuit-to-tensor} e o demo do AlphaTensor-Quantum em pequena escala?}

O projeto será considerado bem-sucedido se, ao final, cumprir simultaneamente os cinco critérios a seguir:
\begin{enumerate}
  \item o demo público do AlphaTensor-Quantum rodar sem modificação estrutural em seus alvos originais;
  \item o pipeline \texttt{compile/resynth/verify} do \texttt{circuit-to-tensor} funcionar ponta a ponta em pelo menos um circuito de controle;
  \item a família de métricas proposta for calculada para todos os circuitos da amostra final;
  \item pelo menos um experimento de RL adaptado na Faixa S for concluído com exportação de decomposição;
  \item a tabela final incluir entre $8$ e $15$ circuitos viáveis, ou, caso a disponibilidade pública imponha um conjunto menor, todos os circuitos viáveis acompanhados de justificativa técnica.
\end{enumerate}

\section{Scope and Non-Scope}

\subsection*{Escopo obrigatório}

O escopo obrigatório do projeto é fechado e executável:
\begin{itemize}
  \item reproduzir apenas o que é publicamente e realisticamente reexecutável;
  \item rodar o demo público do \texttt{alphatensor\_quantum};
  \item montar e validar o pipeline \texttt{compile/resynth/verify} via \texttt{circuit-to-tensor};
  \item rodar baselines com PyZX;
  \item aproveitar decomposições públicas do diretório \texttt{/decompositions} sempre que houver correspondência com circuitos do benchmark;
  \item implementar e medir uma \textbf{família} de métricas de estrutura Clifford--não-Clifford;
  \item produzir comparação quantitativa entre: circuito original, PyZX, QuiZX+\texttt{compile}, demo RL e decomposições públicas quando disponíveis.
\end{itemize}

\subsection*{Não-escopo obrigatório}

Os itens abaixo ficam explicitamente fora do escopo:
\begin{itemize}
  \item reproduzir integralmente todos os números do artigo-base;
  \item treinar agentes gerais no regime computacional do paper;
  \item cobrir de maneira homogênea todos os circuitos grandes do benchmark \textit{op-T-mize};
  \item reconstruir infraestrutura interna da DeepMind;
  \item escrever o relatório como se o repositório público fosse equivalente ao stack interno usado no paper;
  \item depender de qualquer resultado privado, não reprodutível ou inacessível publicamente.
\end{itemize}

\subsection*{Partição operacional do benchmark}

O benchmark será particionado depois da etapa de compilação tensorial, usando o tamanho do tensor compilado como critério operacional:

\begin{itemize}
  \item \textbf{Faixa S} (core com treino RL permitido): alvos do demo público e circuitos adicionais cujo tensor compilado tenha tamanho $\leq 8$. Nesta faixa é permitido adaptar ou treinar o demo.
  \item \textbf{Faixa M} (core sem RL novo obrigatório): circuitos com tensores compilados de tamanho entre $9$ e $20$. Nesta faixa o foco será \texttt{compile}, \texttt{resynth}, PyZX, QuiZX, decomposições públicas e métricas estruturais.
  \item \textbf{Faixa L} (fora do core deliverable): circuitos com tensores compilados de tamanho $>20$. Esta faixa será tratada apenas como exploração opcional, apêndice ou análise estrutural; ela não é dependência para o sucesso do projeto.
\end{itemize}

\subsection*{Amostra final e regra de seleção}

A amostra final seguirá as regras abaixo:
\begin{enumerate}
  \item compilar todos os circuitos públicos disponíveis do benchmark selecionado;
  \item atribuir cada circuito à Faixa S, M ou L;
  \item incluir \textbf{todos} os circuitos da Faixa S;
  \item incluir até \textbf{10} circuitos da Faixa M, escolhidos por amostragem estratificada em três níveis de T-count inicial (baixo, médio e alto) e, quando possível, cobrindo famílias distintas de aplicações;
  \item não incluir circuitos da Faixa L na avaliação principal, salvo como estudo opcional claramente marcado.
\end{enumerate}

Essa decisão fecha o escopo e impede que o projeto se torne inviável por tentar escalar treino RL para casos que pertencem ao regime computacional do artigo, e não ao regime computacional disponível neste projeto.

\section{Scientific Motivation}

A motivação científica do projeto é dupla.

Primeiro, \textbf{T-gates são caros} em computação quântica tolerante a falhas, de modo que reduzir T-count e T-depth continua sendo um objetivo central em compilação quântica. O AlphaTensor-Quantum é interessante porque conecta esse problema à decomposição tensorial e a deep reinforcement learning, ao invés de tratá-lo apenas como simplificação simbólica local. Isso o torna um excelente ponto de partida para um projeto de Deep Learning que também tenha relevância em otimização quântica.

Segundo, a utilidade de um método de otimização não deve ser medida apenas pelo menor T-count final. Em prática experimental, há três questões adicionais:
\begin{itemize}
  \item \textbf{quando} um método melhora mais que outro;
  \item \textbf{em quais tipos de circuitos} ele tende a funcionar melhor;
  \item \textbf{quais propriedades estruturais} do circuito estão associadas ao ganho obtido.
\end{itemize}

É exatamente nesse ponto que a nova contribuição do projeto se encaixa. Em vez de propor um novo grande modelo, o projeto propõe uma \textbf{camada de análise estrutural}: uma família de métricas que capture a composição Clifford/T do circuito, sua segmentação em blocos e sua organização relativa. O valor científico esperado não está condicionado a uma correlação forte; mesmo uma correlação fraca, nula ou instável é informativa, pois ajuda a delimitar o que o demo público consegue ou não capturar.

Do ponto de vista metodológico, o projeto é forte porque combina três virtudes:
\begin{enumerate}
  \item \textbf{reprodução}: o trabalho parte de código público e tenta reexecutá-lo no limite do que é realisticamente possível;
  \item \textbf{extensão}: a proposta adiciona uma família de métricas novas e um protocolo experimental comparativo;
  \item \textbf{controle}: o escopo é fechado por faixas de tensor, orçamento computacional e critérios explícitos de sucesso.
\end{enumerate}

\section{Public Repositories and Their Roles}

\subsection*{\texttt{google-deepmind/alphatensor\_quantum}}

Este repositório será usado para quatro finalidades concretas:
\begin{itemize}
  \item entender a formulação \textit{TensorGame}/RL do método;
  \item rodar o \texttt{src/demo} como núcleo treinável em pequena escala;
  \item reaproveitar o diretório \texttt{/decompositions} como fonte pública de decomposições já obtidas;
  \item extrair uma configuração pública de referência para batch, número de simulações MCTS, frequência de avaliação e comprimento de treino do demo.
\end{itemize}

A regra operacional do projeto é simples: \textbf{o demo será tratado como o único componente público de treino diretamente executável}, enquanto \texttt{/decompositions} será tratado como replay público de resultados, e não como treino reproduzido pelo projeto.

\subsection*{\texttt{tlaakkonen/circuit-to-tensor}}

Este repositório será o eixo da integração entre circuitos e tensores. Seu papel é:
\begin{itemize}
  \item compilar circuitos Clifford+T em tensores binários simétricos;
  \item ressintetizar circuitos a partir de decomposições;
  \item verificar equivalência funcional quando a verificação for viável;
  \item permitir pré-otimização com QuiZX;
  \item servir como ponto de integração entre QASM, tensores, decomposições públicas e circuitos finais.
\end{itemize}

No projeto, \texttt{compile} e \texttt{resynth} são obrigatórios; \texttt{verify} é obrigatório apenas sob orçamento temporal e com possibilidade de status inconclusivo documentado.

\subsection*{\texttt{zxcalc/pyzx}}

PyZX será o \textbf{baseline público principal} de simplificação simbólica. O repositório será usado para:
\begin{itemize}
  \item ler circuitos em QASM;
  \item construir ZX-diagrams;
  \item executar simplificação e extração de circuitos;
  \item gerar um baseline competitivo sem RL;
  \item exportar circuitos otimizados e, quando conveniente, figuras ou artefatos visuais para o relatório.
\end{itemize}

PyZX não será usado apenas como ferramenta de simplificação; ele também servirá de apoio para extração de estatísticas estruturais e inspeção manual de casos qualitativos.

\subsection*{\textit{op-T-mize} da PennyLane}

O benchmark principal do projeto será o \textit{op-T-mize}. Seu papel é:
\begin{itemize}
  \item fornecer um conjunto público de circuitos de referência para T-count optimization;
  \item garantir diversidade suficiente de tamanho e família de circuitos;
  \item permitir comparação com uma base já reconhecida em trabalhos públicos da área;
  \item permitir filtragem por faixas operacionais compatíveis com o orçamento do projeto.
\end{itemize}

A amostra final será derivada do \textit{op-T-mize}, mas o projeto não depende de cobrir os 31 circuitos de forma homogênea com treino RL.

\section{Theory}

\subsection*{Clifford+T}

No formalismo Clifford+T, as portas de Clifford podem ser geradas por combinações de portas como Hadamard ($H$), phase gate ($S$) e CNOT, enquanto a porta $T$ adiciona a não-Cliffordidade necessária para universalidade aproximada. Em computação tolerante a falhas, Clifford gates são relativamente baratas do ponto de vista lógico, enquanto gates do tipo $T$ costumam ser associadas a custos bem maiores de injeção e destilação de estados mágicos. Por isso, otimizar circuitos Clifford+T costuma significar, em primeiro lugar, reduzir o uso efetivo de $T$-gates.

\subsection*{T-count / T-depth}

Duas medidas são centrais:
\begin{itemize}
  \item \textbf{T-count}: número total de gates $T$ e $T^\dagger$ no circuito;
  \item \textbf{T-depth}: número de camadas sequenciais de gates $T/T^\dagger$ que não podem ser paralelizadas sob o critério de conflito em qubits.
\end{itemize}

T-count mede consumo total de gates caros. T-depth mede o grau de serialização do custo não-Clifford. Um circuito pode ter o mesmo T-count de outro e ainda assim possuir T-depth muito maior, o que torna a análise estrutural importante.

\subsection*{Signature tensors}

O AlphaTensor-Quantum reformula a otimização de T-count em termos de um \textbf{tensor de assinatura} binário e simétrico. Em circuitos já reduzidos a uma parte diagonal do tipo CNOT+T, o circuito pode ser representado por um tensor que codifica a estrutura algébrica associada à parte não-Clifford relevante para a otimização.

Para circuitos gerais, a ideia prática é separar o circuito em uma parte diagonal compatível com a representação tensorial e uma parte não diagonal Clifford que permanece como contexto de mudança de base. Essa separação é uma das razões pelas quais faz sentido estudar uma métrica estrutural Clifford--não-Clifford em paralelo ao ganho de T-count.

\subsection*{Waring decomposition}

Uma decomposição do tensor de assinatura em soma de fatores simétricos sobre o corpo binário pode ser escrita, em forma esquemática, como
\[
\mathcal{T} = \sum_{r=1}^{R} u^{(r)} \otimes u^{(r)} \otimes u^{(r)} \pmod 2.
\]
O número de fatores $R$ está diretamente ligado ao custo não-Clifford da realização correspondente. Assim, \textbf{encontrar uma decomposição curta} equivale a encontrar uma realização com T-count reduzido, ao menos até equivalência por gates de Clifford e eventuais esquemas de gadgetização.

\subsection*{TensorGame / RL / MCTS}

O AlphaTensor-Quantum transforma a busca por decomposições curtas em um jogo de decisão sequencial. O estado contém um tensor residual e histórico suficiente das jogadas anteriores; cada ação seleciona um novo fator candidato; a transição atualiza o tensor residual; e a recompensa penaliza o uso de muitos fatores, favorecendo decomposições curtas.

No demo público, essa lógica é acoplada a uma versão simplificada em máquina única usando MCTS e uma rede de política/valor. Para este projeto, essa distinção é crucial: \textbf{o paper original e o demo público não devem ser confundidos}. O demo é uma prova de conceito pública e treinável em pequena escala; o projeto usará essa propriedade deliberadamente, sem extrapolar claims.

\subsection*{Gadgets}

Gadgetização é o uso de padrões estruturados que trocam uma implementação direta em Clifford+T por outra com menor custo equivalente de T-gates, frequentemente ao preço de ancillas, medições ou circuitos auxiliares. No ecossistema deste projeto, gadgets aparecem em dois lugares:
\begin{itemize}
  \item no artigo e no demo, como parte do sinal de recompensa e do problema de otimização;
  \item no \texttt{circuit-to-tensor}, como opção de ressíntese que ativa gadgetização na reconstrução do circuito.
\end{itemize}

Isso exige que o projeto reporte sempre, quando necessário, tanto o circuito produzido quanto o critério de comparação adotado: T-count direto do circuito expandido ou T-count equivalente informado pelo pipeline.

\subsection*{ZX-calculus}

ZX-calculus é uma linguagem gráfica para raciocínio e reescrita de circuitos quânticos via diagramas ZX. Na prática, ela oferece um baseline forte por três motivos:
\begin{enumerate}
  \item representa circuitos em uma forma algébrica flexível;
  \item permite reescritas globais que não são evidentes em uma representação gate-by-gate;
  \item possibilita extrair novamente um circuito simplificado.
\end{enumerate}

PyZX materializa essa abordagem em software público e maduro, tornando-se o baseline natural do projeto.

\subsection*{Motivação para a métrica Clifford--não-Clifford}

T-count sozinho não captura a organização do circuito. Dois circuitos com o mesmo T-count podem diferir muito em:
\begin{itemize}
  \item quantos blocos não-Clifford aparecem;
  \item quão concentrados ou espalhados estão os gates $T$;
  \item quanto contexto Clifford existe entre esses blocos;
  \item quão serializado está o custo não-Clifford.
\end{itemize}

A hipótese de trabalho deste projeto é que uma \textbf{família operacional de métricas estruturais} pode fornecer sinais preditivos sobre:
\begin{itemize}
  \item ganho absoluto de T-count;
  \item ganho percentual;
  \item diferenças entre métodos públicos;
  \item estabilidade estrutural antes e depois da ressíntese.
\end{itemize}

Essa hipótese será testada, não presumida.

\section{Available Compute and Practical Constraints}

O projeto dispõe de \textbf{duas GPUs NVIDIA A100}. Esse recurso é forte para um projeto acadêmico, mas é muito inferior ao regime computacional reportado no artigo-base. Portanto, a política de uso computacional será a seguinte:

\subsection*{Alocação de recursos}

\begin{itemize}
  \item \textbf{GPU A}: prioritariamente dedicada a treino do demo público e pequenas ablações associadas.
  \item \textbf{GPU B}: dedicada a reruns, segunda semente quando realmente necessária, inferência auxiliar, pós-processamento JAX e experimentos controlados de adaptação do demo.
  \item \textbf{CPU}: usada para \texttt{compile}, \texttt{resynth}, parsing, normalização de QASM, PyZX, agregação de resultados, geração de tabelas e scripts auxiliares em Rust e Python.
\end{itemize}

\subsection*{Caps obrigatórios de execução}

Para impedir que o escopo se expanda de forma irrealista, o projeto adotará os seguintes limites:
\begin{itemize}
  \item no máximo \textbf{8 jobs de treino RL} em toda a execução do projeto;
  \item no máximo \textbf{2 sementes} para qualquer configuração RL incluída na tabela principal;
  \item no máximo \textbf{12 horas} de wall-clock por run RL adaptado;
  \item no máximo \textbf{6 horas} de wall-clock por run de validação do demo padrão;
  \item no máximo \textbf{20 minutos} por circuito para PyZX na etapa principal;
  \item no máximo \textbf{30 minutos} por circuito de Faixa S e \textbf{60 minutos} por circuito de Faixa M para \texttt{verify}; acima disso, o status será marcado como \texttt{timeout} ou \texttt{inconclusive};
  \item a Faixa L não terá verificação obrigatória.
\end{itemize}

\subsection*{Configuração inicial do demo}

A configuração inicial do demo seguirá a configuração pública do repositório como ponto de partida: batch size $128$, $80$ simulações MCTS por passo, $50\,000$ passos de treino e avaliação a cada $50$ passos. Essa configuração será tratada como \textbf{baseline operacional pública}, não como hiperparâmetro ótimo universal. Se a execução for instável ou lenta demais, o primeiro recuo permitido será reduzir MCTS de $80$ para $40$, mantendo o restante fixo. Não será feita busca ampla de hiperparâmetros.

\subsection*{Early stopping e fallback}

As regras de early stopping serão:
\begin{itemize}
  \item após os primeiros $10\,000$ passos, encerrar um run RL se o melhor T-count não melhorar por $5\,000$ passos consecutivos;
  \item encerrar imediatamente qualquer run que exceda o cap de wall-clock;
  \item em caso de falha repetida na adaptação RL para um circuito da Faixa S, usar apenas PyZX, QuiZX+\texttt{compile} e, se houver, decomposições públicas.
\end{itemize}

\subsection*{Princípio de honestidade computacional}

O documento final e todos os resultados deverão diferenciar explicitamente:
\begin{itemize}
  \item \textbf{o que o paper fez} em regime de TPU e larga escala;
  \item \textbf{o que este projeto executará} com duas A100 e código público parcial.
\end{itemize}

\section{Execution Plan}

\subsection*{Estrutura sugerida do repositório local}

O projeto deverá ser organizado em um repositório local com a seguinte estrutura:

\begin{verbatim}
project/
  external/
    alphatensor_quantum/
    circuit-to-tensor/
    pyzx/
  configs/
    demo/
    metrics/
    experiments/
  data/
    raw/
      op_t_mize/
      demo_targets/
    normalized_qasm/
    compiled/
    decompositions_public/
    decompositions_demo/
    pyzx_outputs/
    resynth_outputs/
    verification/
  logs/
    compile/
    pyzx/
    demo/
    verify/
  results/
    csv/
    figures/
    tables/
  reports/
  scripts/
\end{verbatim}

\subsection*{Ambientes e dependências}

O plano adota duas pilhas principais:
\begin{itemize}
  \item \textbf{Python 3.11}: demo JAX/MCTX, PyZX, scripts de análise, agregação e geração de tabelas;
  \item \textbf{Rust estável}: \texttt{circuit-to-tensor}.
\end{itemize}

Dependências mínimas recomendadas:
\begin{itemize}
  \item JAX com backend GPU compatível com o ambiente CUDA do cluster;
  \item dependências do \texttt{alphatensor\_quantum/src/demo/requirements.txt};
  \item PyZX em versão estável recente;
  \item PennyLane apenas para ingestão pública do benchmark, se isso facilitar o download;
  \item Rust via \texttt{rustup};
  \item \texttt{feynver} no \texttt{PATH} para verificação quando aplicável.
\end{itemize}

\subsection*{Scripts planejados}

A execução será organizada pelos scripts da Tabela~\ref{tab:scripts}. Os nomes são parte do plano e devem ser mantidos para garantir rastreabilidade.

\begin{table}[ht]
\centering
\caption{Scripts planejados e responsabilidade de cada um.}
\label{tab:scripts}
\begin{tabularx}{\textwidth}{>{\ttfamily}p{0.33\textwidth} Y}
\toprule
Script & Responsabilidade \\
\midrule
scripts/setup\_env.sh & Criar ambiente Python, instalar dependências do demo, instalar PyZX, instalar Rust e compilar \texttt{circuit-to-tensor}. \\
scripts/download\_or\_prepare\_benchmarks.py & Obter os circuitos públicos do benchmark, normalizar nomes, montar catálogo inicial e metadados brutos. \\
scripts/export\_qasm\_v2.py & Converter ou reexportar circuitos para OpenQASM v2 quando necessário. \\
scripts/run\_pyzx\_baseline.py & Ler QASM, simplificar com PyZX, extrair circuito otimizado, salvar QASM final e estatísticas. \\
scripts/run\_compile\_pipeline.sh & Executar \texttt{compile} com e sem QuiZX, armazenar tensores, matrizes, logs e saídas intermediárias. \\
scripts/run\_resynth.sh & Executar \texttt{resynth} a partir de decomposições públicas ou exportadas pelo demo, com e sem gadgetização. \\
scripts/run\_verify.sh & Executar verificação por \texttt{feynver} sob orçamento temporal. \\
scripts/run\_demo\_train.py & Rodar o demo padrão e o demo adaptado na Faixa S. \\
scripts/export\_best\_demo\_decompositions.py & Salvar a melhor sequência de fatores encontrada pelo demo em formato \texttt{.npy}. \\
scripts/compute\_metrics.py & Calcular T-count, T-depth, métricas estruturais e tempos por método e por circuito. \\
scripts/aggregate\_results.py & Unificar resultados em CSV único, validar chaves e produzir tabela canônica por circuito/método. \\
scripts/make\_tables.py & Gerar tabelas finais em CSV e, opcionalmente, em \LaTeX. \\
scripts/plot\_results.py & Gerar gráficos de dispersão, correlações e comparações por método. \\
\bottomrule
\end{tabularx}
\end{table}

\subsection*{Etapa 1: ingestão de benchmarks}

A primeira execução prática será:
\begin{enumerate}
  \item obter os circuitos públicos do \textit{op-T-mize};
  \item garantir representação em OpenQASM v2;
  \item normalizar nomes, número de qubits e estatísticas básicas;
  \item armazenar um catálogo em CSV com colunas mínimas:
  \[
  \texttt{circuit\_id, source, n\_qubits, raw\_gate\_count, qasm\_path.}
  \]
\end{enumerate}

Regra de exclusão: qualquer circuito que, após \textbf{duas tentativas determinísticas de normalização}, ainda não puder ser exportado para OpenQASM v2 ficará fora do core deliverable, mas será listado em apêndice como \texttt{not-normalized}.

\subsection*{Etapa 2: baselines de PyZX}

Para cada circuito em QASM v2:
\begin{enumerate}
  \item carregar no PyZX;
  \item converter para representação de grafo;
  \item aplicar rotina de simplificação;
  \item extrair circuito otimizado;
  \item salvar o QASM otimizado;
  \item calcular T-count, T-depth, tempo de execução e status.
\end{enumerate}

Regra de falha: se a extração do circuito falhar, o circuito não será usado para comparação PyZX, mas continuará elegível para os demais métodos.

\subsection*{Etapa 3: compilação tensorial}

Cada circuito QASM v2 será processado de duas formas:
\begin{enumerate}
  \item \textbf{sem QuiZX}: \texttt{compile} padrão;
  \item \textbf{com QuiZX}: \texttt{compile -z}.
\end{enumerate}

Em ambos os casos, a emissão mínima obrigatória será:
\begin{itemize}
  \item \texttt{tensor};
  \item \texttt{matrix};
  \item \texttt{circuit-qasm};
  \item \texttt{log}.
\end{itemize}

A verificação será executada separadamente por \texttt{scripts/run\_verify.sh}, para permitir orçamentos distintos e evitar que a compilação principal fique bloqueada.

\subsection*{Etapa 4: partição por faixas}

Após a compilação, cada circuito será classificado pela dimensão do tensor compilado:
\begin{itemize}
  \item Faixa S: tamanho $\leq 8$;
  \item Faixa M: tamanho entre $9$ e $20$;
  \item Faixa L: tamanho $>20$.
\end{itemize}

O script \texttt{scripts/aggregate\_results.py} deve gerar uma tabela \texttt{circuit\_inventory.csv} contendo:
\[
\texttt{circuit\_id, n\_qubits, tensor\_size, faixa, compile\_status, pyzx\_status.}
\]

\subsection*{Etapa 5: uso de decomposições públicas}

O diretório \texttt{/decompositions} do repositório \texttt{alphatensor\_quantum} contém, para cada nome de circuito suportado, até dez decomposições não equivalentes em arranjos booleanos. O projeto usará essas decomposições apenas como \textbf{replay público}.

Procedimento:
\begin{enumerate}
  \item mapear nomes de circuito entre benchmark local e arquivo de decomposições;
  \item quando houver correspondência exata, extrair no máximo as \textbf{3 primeiras} decomposições para ressíntese;
  \item executar \texttt{resynth} com e sem gadgetização;
  \item verificar quando o orçamento permitir;
  \item manter o melhor resultado validado ou, na ausência de validação, o melhor resultado com status claramente marcado.
\end{enumerate}

\subsection*{Etapa 6: demo padrão}

O demo padrão deverá ser rodado \textbf{sem adaptação estrutural} para validar a infraestrutura pública. Os alvos padrão serão exatamente os do repositório público. O objetivo desta etapa não é produzir o melhor número do projeto, mas comprovar:
\begin{itemize}
  \item que o ambiente JAX/MCTX funciona;
  \item que o log de treinamento é compreensível;
  \item que o projeto consegue exportar decomposições e conectá-las ao \texttt{resynth}.
\end{itemize}

\subsection*{Etapa 7: demo adaptado na Faixa S}

O demo adaptado será usado apenas na Faixa S. As decisões são:
\begin{itemize}
  \item agrupar circuitos por tamanho tensorial e afinidade temática quando possível;
  \item no máximo \textbf{3 tensores por run} adaptado;
  \item usar a configuração pública como baseline inicial;
  \item realizar no máximo \textbf{2 sementes} por configuração selecionada;
  \item exportar a melhor sequência de fatores assim que um novo melhor retorno for encontrado em episódio terminal.
\end{itemize}

Modificação obrigatória do demo: criar um mecanismo de logging que salve, para o melhor episódio terminal observado, a sequência de fatores em formato \texttt{.npy} compatível com o \texttt{resynth}. Essa modificação faz parte do projeto e deve ser documentada em apêndice metodológico.

\subsection*{Etapa 8: ressíntese}

A ressíntese será aplicada apenas quando existir decomposição:
\begin{itemize}
  \item decomposição pública proveniente de \texttt{/decompositions};
  \item decomposição exportada pelo demo;
  \item eventualmente, decomposição heurística adicional se for implementada dentro do orçamento e sem alterar o escopo central.
\end{itemize}

Para cada decomposição, executar:
\begin{enumerate}
  \item \texttt{resynth} sem gadgetização;
  \item \texttt{resynth} com gadgetização, quando aplicável;
  \item cálculo de métricas e verificação sob cap temporal.
\end{enumerate}

\subsection*{Etapa 9: verificação}

A verificação com \texttt{feynver} será tratada como forte evidência de correção, mas não como requisito universal. Os status permitidos são:
\begin{itemize}
  \item \texttt{verified};
  \item \texttt{timeout};
  \item \texttt{inconclusive};
  \item \texttt{not-run}.
\end{itemize}

A tabela final sempre exibirá o status de verificação, para impedir que circuitos não verificados sejam apresentados como equivalentes sem ressalva.

\subsection*{Etapa 10: cálculo de métricas e agregação}

Todos os circuitos comparados serão convertidos para uma base normalizada de portas, preferencialmente Clifford+T explícita. O arquivo canônico final deverá ter uma linha por par \texttt{(circuito, método)} e as colunas:
\begin{verbatim}
circuit_id, method, source, faixa, n_qubits, tensor_size,
verify_status, runtime_sec,
tcount_before, tcount_after, delta_t, rel_gain_t,
tdepth_before, tdepth_after, delta_tdepth,
rho_t, rho_w, n_clifford_blocks, n_nonclifford_blocks,
avg_nonclifford_block_len, hadamard_boundary_density,
tdepth_over_tcount
\end{verbatim}

\section{Experimental Matrix}

\subsection*{Grupos de circuitos}

\begin{table}[ht]
\centering
\caption{Grupos experimentais e papel de cada grupo.}
\label{tab:groups}
\begin{tabularx}{\textwidth}{p{0.15\textwidth} p{0.22\textwidth} Y Y}
\toprule
Grupo & Regra de inclusão & Métodos obrigatórios & Papel no projeto \\
\midrule
Demo-controle & Alvos originais do demo público & Demo padrão, PyZX, métricas, ressíntese quando houver decomposição exportada & Validar infraestrutura pública e comportamento base do demo \\
Faixa S & Tensor compilado $\leq 8$ & PyZX, QuiZX+\texttt{compile}, demo adaptado, métricas, ressíntese & Núcleo com treino RL permitido \\
Faixa M & Tensor compilado entre $9$ e $20$ & PyZX, QuiZX+\texttt{compile}, decomposições públicas quando existirem, métricas & Núcleo comparativo sem depender de novo treino RL \\
Faixa L & Tensor compilado $>20$ & Métricas e análise opcional; demais métodos apenas se o custo for baixo & Apêndice e delimitação de limite prático \\
\bottomrule
\end{tabularx}
\end{table}

\subsection*{Métodos comparados}

\begin{table}[ht]
\centering
\caption{Métodos da matriz comparativa.}
\label{tab:methods}
\begin{tabularx}{\textwidth}{p{0.18\textwidth} p{0.12\textwidth} Y Y}
\toprule
Método & Treino novo & Saída principal & Papel analítico \\
\midrule
Original & Não & Circuito de referência & Ponto de partida para T-count, T-depth e métricas estruturais \\
PyZX & Não & Circuito simplificado por ZX-calculus & Baseline público forte sem RL \\
QuiZX+\texttt{compile} & Não & Circuito pré-otimizado e tensores compilados & Baseline tensorial sem treino novo \\
Decomposição pública+\texttt{resynth} & Não & Circuito ressintetizado a partir de replay público & Comparação com resultados públicos do ecossistema AlphaTensor-Quantum \\
Demo padrão & Sim, mas apenas no alvo público & Logs e decomposição exportada & Validação do componente treinável público \\
Demo adaptado & Sim, restrito à Faixa S & Decomposição exportada e circuito ressintetizado & Estudo RL em pequena escala compatível com o projeto \\
\bottomrule
\end{tabularx}
\end{table}



\section{Metrics}

\subsection*{Métricas de ganho}

As métricas primárias de desempenho serão:
\[
\Delta T = T_{\text{antes}} - T_{\text{depois}},
\qquad
g_T = \frac{T_{\text{antes}} - T_{\text{depois}}}{\max(T_{\text{antes}},1)}.
\]
De forma análoga, serão calculadas:
\[
\Delta D_T = D_{T,\text{antes}} - D_{T,\text{depois}},
\]
onde $D_T$ é a T-depth.

\subsection*{Métrica 1: fração não-Clifford}

A primeira camada da família de métricas é a fração não-Clifford:
\[
\rho_T = \frac{N_T}{N_{\text{total}}},
\]
onde $N_T$ conta portas $T$ e $T^\dagger$, e $N_{\text{total}}$ conta todas as portas na forma normalizada de comparação.

A variante ponderada por custo será
\[
\rho_w^{(\lambda)} = \frac{\lambda N_T}{N_C + \lambda N_T},
\]
onde $N_C$ é o número de portas Clifford na representação normalizada. A análise principal usará $\lambda = 10$, enquanto a ablação testará $\lambda \in \{5,10,20\}$.

Essa escolha é deliberada: evita assumir um custo físico exato dependente de arquitetura e, ao mesmo tempo, preserva a ideia de que T-gates são significativamente mais caros que gates Clifford.

\subsection*{Métrica 2: estrutura de separação}

A segunda camada da família mede a estrutura de segmentação. Após normalização do circuito para uma sequência linear de portas em base Clifford+T:
\begin{itemize}
  \item define-se \textbf{bloco Clifford} como um segmento maximal de portas do conjunto Clifford;
  \item define-se \textbf{bloco não-Clifford} como um segmento maximal de portas do conjunto $\{T,T^\dagger\}$.
\end{itemize}

As métricas estruturais obrigatórias serão:
\begin{itemize}
  \item $B_C$: número de blocos Clifford;
  \item $B_{NC}$: número de blocos não-Clifford;
  \item $\bar{\ell}_{NC}$: comprimento médio dos blocos não-Clifford;
  \item $H_{\partial}$: densidade de Hadamard nas fronteiras entre blocos Clifford e não-Clifford;
  \item T-count;
  \item T-depth;
  \item razão
  \[
  \eta_T = \frac{D_T}{\max(T,1)}.
  \]
\end{itemize}

A densidade de Hadamard na fronteira será definida por
\[
H_{\partial} = \frac{N_H^{(\partial)}}{\max(B_C + B_{NC} - 1,1)},
\]
onde $N_H^{(\partial)}$ conta gates $H$ a uma distância de até uma porta de uma fronteira entre bloco Clifford e bloco não-Clifford. A janela de fronteira $w=1$ será a definição principal; a ablação testará $w=2$.

\subsection*{Normalização de portas}

Para cálculo das métricas, todos os circuitos serão convertidos, sempre que possível, para uma base comparável que explicite:
\[
\{H, S, S^\dagger, X, Y, Z, \mathrm{CNOT}, T, T^\dagger\}.
\]
Se um método produzir portas compostas, a regra será:
\begin{enumerate}
  \item tentar expandir para Clifford+T canônico;
  \item se a expansão não for suportada pelo pipeline público, registrar o circuito apenas com métricas disponíveis e marcar a ausência de comparabilidade estrutural total.
\end{enumerate}

\subsection*{Perguntas experimentais obrigatórias}

A avaliação quantitativa deverá responder às cinco perguntas abaixo:
\begin{enumerate}
  \item A métrica prevê ganho absoluto de T-count?
  \item A métrica prevê ganho percentual?
  \item PyZX e AlphaTensor-Quantum melhoram os mesmos circuitos?
  \item Há circuitos em que o baseline ZX vence o demo RL?
  \item A métrica é estável antes e depois da ressíntese?
\end{enumerate}

\subsection*{Procedimento estatístico}

O procedimento estatístico será fechado e simples:
\begin{itemize}
  \item correlação de \textbf{Spearman} como análise primária entre cada métrica estrutural e os ganhos $\Delta T$ e $g_T$;
  \item correlação de \textbf{Pearson} como análise secundária;
  \item intervalos de confiança por bootstrap com $1000$ reamostragens;
  \item gráficos de dispersão por método para inspeção qualitativa;
  \item sem ``p-hacking'': o foco será tamanho de efeito, estabilidade do sinal e interpretação metodológica.
\end{itemize}

\section{Ablations}

As ablações serão pequenas, controladas e limitadas pelo orçamento computacional. O projeto incluirá obrigatoriamente as seguintes ablações:

\begin{enumerate}
  \item \textbf{Com e sem gadgets}: executar demo padrão e, quando possível, demo adaptado e \texttt{resynth} com gadgetização ligada e desligada.
  \item \textbf{Com e sem pré-otimização QuiZX}: para todos os circuitos da amostra final, comparar \texttt{compile} padrão versus \texttt{compile -z}.
  \item \textbf{Com diferentes definições da métrica}: variar $\lambda \in \{5,10,20\}$ em $\rho_w^{(\lambda)}$ e usar janela de fronteira $w \in \{1,2\}$ em $H_{\partial}$.
  \item \textbf{Com e sem decomposições públicas}: em circuitos com correspondência pública, comparar a tabela final incluindo ou removendo o replay de \texttt{/decompositions}.
  \item \textbf{Demo padrão versus demo levemente adaptado}: comparar o comportamento do demo nos alvos originais e em um conjunto pequeno da Faixa S.
\end{enumerate}

Escopo das ablações:
\begin{itemize}
  \item as ablações de gadgetização e demo adaptado serão restritas a no máximo \textbf{2 circuitos de controle} e \textbf{2 circuitos da Faixa S};
  \item a ablação de QuiZX será aplicada em toda a amostra final, pois não exige treino novo;
  \item a ablação de métrica será aplicada a todos os circuitos da amostra final, pois é computacionalmente barata.
\end{itemize}

\section{Risks and Mitigations}

\begin{longtable}{p{0.21\textwidth} p{0.18\textwidth} p{0.26\textwidth} p{0.25\textwidth}}
\caption{Riscos principais e mitigação.}
\label{tab:risks}\\
\toprule
Risco & Impacto & Mitigação & Fallback \\
\midrule
\endfirsthead
\toprule
Risco & Impacto & Mitigação & Fallback \\
\midrule
\endhead
Demo não generaliza bem & O RL pequeno não melhora circuitos novos de Faixa S & Limitar claims; tratar RL como prova de conceito; manter demo padrão como controle obrigatório & Basear conclusão comparativa em PyZX, QuiZX+\texttt{compile} e decomposições públicas \\
Conversão para tensor falha & Parte do benchmark não entra no pipeline tensorial & Exigir OpenQASM v2; normalização determinística; registrar exclusões com razão técnica & Trabalhar apenas com subconjunto viável e documentar cobertura efetiva \\
Verificação lenta ou inconclusiva & Dificuldade em afirmar equivalência formal de todos os casos & Caps de tempo; status explícitos \texttt{verified}/\texttt{timeout}/\texttt{inconclusive} & Relatar equivalência apenas quando verificada; manter demais casos como candidatos não verificados \\
Métrica não correlaciona com ganho & A principal hipótese exploratória pode falhar & Definir família de métricas, não uma única; usar análise estatística transparente & Tratar resultado negativo como achado metodológico legítimo \\
Custo computacional alto & RL e verificação consumirem o cronograma & Wall-clock caps, early stopping, limite de 8 runs RL, foco em Faixa S/M & Encerrar RL adaptado cedo e reforçar análise com replay público \\
Poucos circuitos entram na Faixa S & RL adaptado perde massa experimental & Manter demo-controle e usar toda Faixa S viável, ainda que pequena & Redefinir RL como estudo de caso e concentrar comparação ampla na Faixa M \\
PyZX falha na extração em parte dos casos & Baseline incompleto & Separar falha de leitura, simplificação e extração; manter logs & Comparar apenas nos circuitos onde PyZX produz circuito final \\
Mapeamento de nomes nas decomposições públicas é ambíguo & Risco de replay incorreto & Exigir nome exato ou checagem adicional por dimensão e estatísticas do tensor & Se não houver correspondência inequívoca, não usar aquela decomposição \\
\bottomrule
\end{longtable}

\section{Deliverables}

Ao final do projeto, os entregáveis obrigatórios serão:

\begin{enumerate}
  \item \textbf{código organizado} em repositório único, com scripts nomeados conforme a Tabela~\ref{tab:scripts};
  \item \textbf{arquivo de ambiente} (\texttt{requirements.txt}, \texttt{environment.yml} ou equivalente) e instruções de reprodução;
  \item \textbf{CSVs de métricas} com schema canônico por circuito e por método;
  \item \textbf{tabelas finais} em formato CSV e, opcionalmente, \LaTeX;
  \item \textbf{gráficos} de comparação por método e correlação métrica--ganho;
  \item \textbf{apêndice metodológico} descrevendo adaptações do demo, normalização de QASM e política de verificação;
  \item \textbf{relatório final} com resultados, limitações e interpretação;
  \item \textbf{checklist de reprodutibilidade} contendo:
  \begin{itemize}
    \item hashes de commit dos repositórios externos;
    \item versões de Python, JAX, PyZX e Rust;
    \item flags usadas em \texttt{compile}, \texttt{resynth} e demo;
    \item orçamento temporal real por bloco;
    \item status de verificação de cada circuito final.
  \end{itemize}
\end{enumerate}

Arquivos mínimos esperados no diretório \texttt{results/}:
\begin{itemize}
  \item \texttt{results/csv/circuit\_inventory.csv};
  \item \texttt{results/csv/final\_metrics.csv};
  \item \texttt{results/csv/final\_summary.csv};
  \item \texttt{results/tables/final\_comparison.tex};
  \item \texttt{results/figures/*.pdf}.
\end{itemize}


Se qualquer semana crítica atrasar, a prioridade de preservação será:
\[
\text{PyZX + QuiZX+\texttt{compile} + métricas + demo-controle} \;>\; \text{demo adaptado} \;>\; \text{Faixa L}.
\]

\section{Expected Outcomes}

Os resultados esperados devem ser formulados em termos de \textbf{viabilidade, comparação e delimitação}, e não como promessas de superar o paper.

Espera-se obter:
\begin{itemize}
  \item um pipeline público funcional que conecte OpenQASM v2, \texttt{circuit-to-tensor}, PyZX, demo público do AlphaTensor-Quantum e análise quantitativa;
  \item uma comparação metodológica clara entre simplificação por ZX-calculus, pré-otimização QuiZX+\texttt{compile} e uso do ecossistema AlphaTensor-Quantum em pequena escala;
  \item uma avaliação objetiva da utilidade da família de métricas estruturais proposta;
  \item uma delimitação prática dos limites do demo público sob orçamento compatível com duas A100;
  \item um conjunto de resultados com valor científico mesmo se a correlação entre métricas e ganho for fraca, instável ou negativa.
\end{itemize}

Em particular, o projeto deverá ser capaz de produzir uma das seguintes conclusões válidas:
\begin{enumerate}
  \item a família de métricas fornece sinal preditivo útil;
  \item a família de métricas fornece sinal fraco, mas interpretável apenas em certos regimes;
  \item a família de métricas não explica adequadamente o ganho, o que constitui um resultado negativo relevante e bem documentado.
\end{enumerate}

\section{References}

\begin{thebibliography}{99}

\bibitem{ruiz2025nature}
F. J. R. Ruiz, T. Laakkonen, J. Bausch, M. Balog, M. Barekatain, F. J. H. Heras, A. Novikov, N. Fitzpatrick, B. Romera-Paredes, J. van de Wetering, A. Fawzi, K. Meichanetzidis e P. Kohli.
\newblock \emph{Quantum Circuit Optimization with AlphaTensor}.
\newblock Nature Machine Intelligence, 2025.
\newblock Disponível em: \url{https://www.nature.com/articles/s42256-025-01001-1}.

\bibitem{ruiz2024arxiv}
F. J. R. Ruiz, T. Laakkonen, J. Bausch, M. Balog, M. Barekatain, F. J. H. Heras, A. Novikov, N. Fitzpatrick, B. Romera-Paredes, J. van de Wetering, A. Fawzi, K. Meichanetzidis e P. Kohli.
\newblock \emph{Quantum Circuit Optimization with AlphaTensor}.
\newblock arXiv:2402.14396, 2024.
\newblock Disponível em: \url{https://arxiv.org/abs/2402.14396}.

\bibitem{alphatensorrepo}
Google DeepMind.
\newblock \emph{alphatensor\_quantum}.
\newblock Repositório público no GitHub.
\newblock Disponível em: \url{https://github.com/google-deepmind/alphatensor_quantum}.

\bibitem{circuittotensorrepo}
T. Laakkonen.
\newblock \emph{circuit-to-tensor}.
\newblock Repositório público no GitHub.
\newblock Disponível em: \url{https://github.com/tlaakkonen/circuit-to-tensor}.

\bibitem{pyzxrepo}
A. Kissinger, J. van de Wetering e colaboradores.
\newblock \emph{PyZX}.
\newblock Repositório público no GitHub.
\newblock Disponível em: \url{https://github.com/zxcalc/pyzx}.

\bibitem{kissinger2020tcount}
A. Kissinger e J. van de Wetering.
\newblock \emph{Reducing T-count with the ZX-calculus}.
\newblock Physical Review A 102, 022406, 2020.
\newblock Pré-print disponível em: \url{https://arxiv.org/abs/1903.10477}.

\bibitem{kissinger2019pyzx}
A. Kissinger e J. van de Wetering.
\newblock \emph{PyZX: Large Scale Automated Diagrammatic Reasoning}.
\newblock arXiv:1904.04735, 2019.
\newblock Disponível em: \url{https://arxiv.org/abs/1904.04735}.

\bibitem{optmize}
PennyLane.
\newblock \emph{op-T-mize dataset}.
\newblock Página pública do benchmark.
\newblock Disponível em: \url{https://pennylane.ai/datasets/op-t-mize}.

\bibitem{deepmindpub}
Google DeepMind.
\newblock \emph{AlphaTensor for Optimizing Quantum Computations}.
\newblock Página pública da publicação.
\newblock Disponível em: \url{https://deepmind.google/research/publications/77240/}.

\end{thebibliography}

\end{document}